\chapter{INTRODUÇÃO}

Este capítulo contextualiza o problema que motiva o trabalho, apresentando o cenário atual de ameaças cibernéticas e o papel crescente da inteligência artificial na segurança de redes. A partir desse contexto, são formuladas as hipóteses de pesquisa, definidos os objetivos do estudo e descrita a metodologia adotada para a condução dos experimentos comparativos entre os modelos \textit{Random Forest} e \textit{Deep Neural Network}.

\section{Contextualização}

A elaboração de estratégias de defesa cibernética robustas tornou-se uma exigência crítica e inadiável no cenário tecnológico contemporâneo. A digitalização acelerada dos processos corporativos, governamentais e sociais, impulsionada pela adoção massiva de computação em nuvem e dispositivos móveis, ampliou a superfície de ataque e introduziu vulnerabilidades que, quando exploradas, geram impactos econômicos, operacionais e reputacionais severos. Nesse contexto, preservar a confidencialidade, a integridade e a disponibilidade da informação tornou-se um desafio dinâmico que exige evolução constante frente a adversários cada vez mais organizados.

A análise dos relatórios recentes revela uma mudança de tendência nos prejuízos financeiros decorrentes de ciberataques. Estes dados indicam que as defesas tradicionais, baseadas puramente em regras estáticas e assinaturas, já não são suficientes para conter a sofisticação das ameaças modernas, resultando em custos globais de violação de dados que atingiram patamares históricos \parencite[4]{IBM2024}.

Segundo o relatório \textit{Cost of a Data Breach Report 2025}, o custo médio global de uma violação de dados apresentou uma queda para US\$4,44 milhões, representando uma redução de 9\% em relação ao ano anterior \parencite[3]{IBM2025}. Essa diminuição é atribuída diretamente à redução do ciclo de vida das violações: o tempo médio para identificar e conter um incidente caiu para 241 dias (o menor patamar em nove anos), impulsionado fundamentalmente pelo uso extensivo de Inteligência Artificial e automação pelas equipes de segurança \parencite[7]{IBM2025}.

No Brasil, o cenário acompanhou a tendência global de redução, com o custo médio de uma violação caindo para US\$1,22 milhão em 2025 \parencite[6]{IBM2025}. Contudo, a complexidade dos ataques permanece crítica. O vetor de ataque inicial mais comum passou a ser o \textit{phishing}, responsável por 16\% das violações, desbancando o uso de credenciais roubadas \parencite[16]{IBM2025}. Além disso, ataques envolvendo comprometimento da cadeia de suprimentos (\textit{supply chain}) tornaram-se o segundo vetor mais frequente e o que demanda maior tempo para identificação e contenção, com o tempo médio de 267 dias \parencite[18]{IBM2025}.

Diante da falibilidade humana e da persistente escassez de profissionais qualificados, que ainda afeta 48\% das organizações \parencite[44]{IBM2025}, a implementação de Sistemas de Detecção de Intrusão (\textit{Intrusion Detection System} - IDS) autônomos baseados em Inteligência Artificial (IA) consolida-se não apenas como uma vantagem, mas como uma necessidade estratégica. Organizações que utilizam \textit{Machine Learning} e \textit{Deep Learning} de forma extensiva economizam, em média, US\$1,9 milhão em custos de violação em relação àquelas que não as adotam, evidenciando a eficácia dessas tecnologias na correlação de eventos complexos e na identificação de anomalias \parencite[4]{IBM2025}.

Dessa forma, escolhemos essas duas tecnologias para fazer um comparativo ao longo desse trabalho para esse laboratório.

\section{Hipóteses}

A pesquisa parte da premissa de que é esperado que algoritmos de \textit{Deep Learning}, em virtude de sua arquitetura baseada em camadas ocultas e capacidade de extração automática de características, apresentem taxas de detecção de intrusão superiores.

Em contrapartida, estima-se que o algoritmo \textit{Random Forest}, pela sua natureza de conjunto de árvores de decisão, oferecerá um desempenho computacional mais eficiente, caracterizado por menores tempos de treinamento e classificação. Essa característica, aliada a uma maior interpretabilidade dos resultados, tende a torná-lo uma solução mais viável para ambientes com recursos de \textit{hardware} limitados, ainda que possa apresentar uma ligeira perda de precisão na identificação de ameaças desconhecidas em comparação às redes neurais profundas.

\section{Objetivos}

O objetivo deste trabalho consiste em desenvolver e implementar um laboratório experimental de detecção de intrusão para realizar uma análise comparativa de desempenho entre técnicas de \textit{Machine Learning} (\textit{Random Forest}) e \textit{Deep Learning} (Redes Neurais Profundas), visando identificar qual abordagem oferece a melhor relação entre eficácia na detecção de ameaças e custo computacional em redes de computadores.

O laboratório é de natureza estritamente experimental e \textit{offline}: os modelos são treinados e avaliados sobre \textit{datasets} de tráfego de rede coletados previamente em ambientes controlados. Não há captura de tráfego em tempo real, conexão com redes corporativas em operação ou implantação operacional dos modelos. O escopo deste trabalho limita-se à comparação sistemática de metodologias de aprendizado de máquina sob condições reproduzíveis, não ao desenvolvimento de um IDS pronto para produção.

Para alcançar essa finalidade, o estudo buscou fundamentar teoricamente os conceitos de Sistemas de Detecção de Intrusão (\textit{IDS}), abordando as diferenças entre detecção por assinatura e por anomalia. Na sequência, foram preparados ambientes de dados a partir de dois \textit{datasets} de referência, o \textit{CICIDS2017} e o \textit{LycoS-IDS2017}, submetidos a pré-processamento, limpeza e normalização. A etapa prática envolveu a implementação, treinamento e validação dos modelos em múltiplos cenários controlados, incluindo experimentos com redução de dimensionalidade por \textit{Mean Decrease in Impurity} (MDI) e uma abordagem não supervisionada via Autoencoder, culminando na avaliação quantitativa dos resultados através de métricas como precisão, revocação, Macro F1-score e análise da matriz de confusão.

\section{Metodologia}

Este trabalho caracteriza-se como uma pesquisa aplicada, de natureza quantitativa e experimental. A metodologia dividiu-se em quatro etapas sequenciais: Levantamento Bibliográfico, Preparação dos Dados, Implementação Experimental e Análise de Resultados.

\subsection{Levantamento Bibliográfico}

A primeira etapa consistiu na revisão sistemática da literatura em bases de dados acadêmicas (Google Scholar, IEEE Xplore, SBC), buscando trabalhos relacionados à aplicação de IA em segurança da informação, com foco em publicações dos últimos cinco anos. Foram consultadas normas técnicas e relatórios de mercado (como IBM Security e Cisco) para contextualizar o cenário de ameaças.

\subsection{Seleção e Pré-processamento do Dataset}

Foram utilizados dois \textit{datasets} públicos de referência para intrusão de rede: o \textit{CICIDS2017} (disponibilizado pelo Canadian Institute for Cybersecurity) e o \textit{LycoS-IDS2017}, uma versão revisada e recurada do \textit{CICIDS2017} com correções de rotulagem e remoção de artefatos de coleta. Ambos contêm tráfego benigno e classes de ataque como \textit{Brute Force}, \textit{DoS}, \textit{Web Attack}, \textit{Infiltration}, \textit{Botnet} e \textit{DDoS}. Os dados passaram por um processo de tratamento (\textit{ETL}), que incluiu:

\begin{itemize}
\item Remoção de valores nulos ou infinitos.
\item Normalização dos atributos numéricos via \textit{StandardScaler} para garantir a convergência eficiente dos algoritmos, especialmente da Rede Neural.
\item Divisão estratificada 70/30 entre treino e teste, preservando a proporção de classes.
\item Redução de dimensionalidade por \textit{Mean Decrease in Impurity} (MDI), aplicada em experimentos adicionais para avaliar o impacto da seleção de atributos sobre o desempenho dos modelos.
\end{itemize}

\subsection{Implementação e Experimentação}

O laboratório virtual foi desenvolvido utilizando a linguagem de programação \textit{Python}. Para o modelo de \textit{Machine Learning}, foi utilizada a biblioteca \textit{Scikit-learn} para implementar o algoritmo \textit{Random Forest}. Para o modelo de \textit{Deep Learning} supervisionado, foram utilizadas as bibliotecas \textit{TensorFlow} e \textit{Keras} para a construção de uma Rede Neural Profunda (\textit{DNN}) com múltiplas camadas densas. Como extensão não supervisionada, foi implementado um Autoencoder treinado exclusivamente com tráfego benigno, com o objetivo de avaliar a viabilidade de detecção de ataques sem dependência de rótulos. O ambiente de execução foi configurado em notebooks (\textit{Jupyter} e \textit{Google Colab}), permitindo a reprodutibilidade dos experimentos.

\subsection{Análise de Dados e Métricas}

A comparação entre os modelos não se baseou apenas na acurácia global, visto que \textit{datasets} de segurança tendem a ser desbalanceados. A análise crítica focou na matriz de confusão e nas métricas de precisão, \textit{recall} e F1-score por classe de ataque, com ênfase no Macro F1-score como indicador de desempenho agregado, observando especialmente a taxa de falsos negativos (ataques não detectados) e falsos positivos (alarmes falsos). Essa abordagem permitiu concluir qual modelo é mais robusto para cada tipo específico de ameaça e sob quais condições de dados.
